\begin{abstract}
Gradient-based meta-learners such as MAML adapt a shared initialization $\theta_0$ to a new task from only a handful of support examples $S_i$, but the mechanism by which $S_i$ steers this adaptation is opaque: when the support set is unrepresentative, hard, or out-of-distribution, adaptation can actively \emph{hurt} query performance, a failure mode known as negative transfer that prior work detects only through accuracy drops, never explains at the feature level. We introduce \textbf{adaptation gain} $\Delta M$, a differentiable, sign-aware measure of how much letting the backbone adapt helps or hurts relative to a feature-reuse baseline where only the classifier head adapts, and \textbf{\fama} (Feature-space Adjoint Meta-learning Attribution), a post-hoc explainer that traces $\Delta M$ back to the support set. \fama extends Grad-CAM across an entire multi-step adaptation trajectory using an adjoint (Pearlmutter Hessian-vector-product) recursion, at $\mathcal{O}(TP)$ cost instead of the $\mathcal{O}(TP^2)$ of naive backpropagation-through-time. We further contribute two faithfulness protocols -- a bidirectional deletion test and cascading-randomization/support-intervention sanity checks -- built for trajectories rather than static predictions. On MAML/Conv4, \fama passes both directions of the deletion test on tieredImageNet ($p<0.001$) and collapses under prior randomization exactly at the layer where the useful gradient signal vanishes; on 600 miniImageNet/tieredImageNet$\to$CUB-200-2011 cross-domain tasks, case-study inspection further points to \emph{contextual bias} -- reliance on background rather than object features -- as a recurring driver of negative transfer.
\end{abstract}
